\documentclass[10pt,a4paper]{amsart}
\usepackage[margin=19mm]{geometry}
\usepackage{amsmath,amssymb,mathtools}
\usepackage[scr=rsfs,cal=euler]{mathalfa}
\usepackage{xfrac}
\usepackage[unicode,hidelinks,bookmarks=false]{hyperref}
\newcommand{\ie}{i.e.\ }
\newcommand{\cf}{cf.\ }
\newcommand{\resp}{respectively\ }
\newcommand{\Subsection}{\subsection}
\providecommand{\nobreakdash}{\nobreak\hbox{-}}
\setlength{\emergencystretch}{2em}
\multlinegap0pt
\usepackage{amssymb}
\usepackage{algorithm}\usepackage{algpseudocode}
\usepackage{mathtools}
\usepackage{tikz}
\usepackage{xcolor}
\newtheorem{assumption}{Assumption}
\newtheorem{lemma}{Lemma}
\newcommand{\QQ}{\mathbb{Q}}
\newcommand{\R}{\mathbb{R}}
\newcommand{\RR}{\mathbb{R}}
\newcommand{\ZZ}{\mathbb{Z}}
\newcommand{\bbq}{\QQ}
\newcommand{\bbz}{\ZZ}
\newcommand{\cE}{\mathcal{E}}
\newcommand{\eps}{\varepsilon}
\newcommand{\problem}[1]{{\color{red}  {Problems: #1}}}
\title{Clustering of large deviations in heavy-tailed moving averages: the catastrophe principle in the long-memory case}
\author{Jiaqi Wang, Gennady Samorodnitsky$^*$}
\date{Source: arXiv:2607.22878v1 (CC BY 4.0)}
\begin{document}
\maketitle

\noindent\emph{Source and licence.} Clustering of large deviations in heavy-tailed moving averages: the catastrophe principle in the long-memory case. Version arXiv:2607.22878v1 (CC BY 4.0), distributed by arXiv under the Creative Commons Attribution 4.0 International licence, http://creativecommons.org/licenses/by/4.0/, as recorded in the arXiv licence metadata for this exact version and shown on the abstract page https://arxiv.org/abs/2607.22878v1. Full licence notice: \texttt{licenses/arxiv-2607.22878-CC-BY-4.0.txt}.\par\smallskip

\noindent\textit{Editorial scope.} Counted unit: the article's lemma lemma:cluster\_equi\_pt (source0.tex lines 1890--1913). The article's complete original proof of it is delivered with it (source0.tex lines 1916--2041, one \texttt{proof} environment of 126 lines). No mathematics is written, repaired or re-derived: the statement and the proof are the article's own lines, in the article's own notation, and no retained line is substituted or re-flowed.  Retained uncounted as the source-local context the proof needs: lines 565--573; lines 988--995; lines 1034--1047; lines 1302--1306.  Nothing was omitted from any retained span, so the package carries no omission record.  No retained line contains a citation, so the wrapper carries no bibliography and no bibliography entry.  No retained line contains a figure environment, an image-inclusion command or drawing code, so no image is delivered and the package declares no asset dependency.  Editorial formatting only, in addition: an AMS \texttt{amsart} preamble that restates the article's own declarations and macros used by the retained lines, namely the environments \texttt{assumption}, \texttt{lemma} on separate counters, together with the standard packages those lines need.  The retained environments print in this document as Assumption 1, Lemma 1 in order of appearance, whereas the article prints the same items with the numbering of its own full document.  The complete native source of the article as distributed in the arXiv e-print for this exact version is bundled unchanged as \texttt{sources/205976/source0.tex}.
\begin{assumption}\label{ass:tail_measure}
$\Gamma \subset \RR^d \setminus
B_{\delta_0}(0)$ for some $\delta_0>0$,  $(pB_++(1-p)B_-)^{-1}\Gamma$ is a
$\nu$-continuity set for $p=0,1$ and almost every $0< p< 1$ (with respect to
the Lebesgue measure) and
$$
\int_\R \nu ( B_y^{-1}\Gamma)\, dy>0. 
$$
\end{assumption}

\begin{equation}\label{eq:coeff_group_err}
\begin{aligned}
\limsup_{n\to\infty} \|\tilde{A}_k^{(n)}/h^{(n)}\| \leq&  \big|(|k| \eps)^{ 1 + \beta}  -
(|k+1|\eps)^{1 + \beta} \big| \\
+& \big||1 - k\eps|^{1 + \beta} - |1 - (k+1)
\eps|^{ 1 + \beta}\big|\max ( \|B _+\|,\|B_-\|), 
\end{aligned}
\end{equation}

\begin{lemma}\label{lemma:asym_lim_meas}
Let $\Phi\subset \RR^d \setminus B_{\delta_0}(0)$. For every $k\in\bbz$, $\eps > 0$, $0< \delta<\delta_0$, 
 $$\limsup_{n\to \infty} \frac{  P({H_{[kn\eps]}^{(n)}} Z_k \in \lambda_n
   \Psi, G_{[kn\eps]}^{(n)} Z_k \in \lambda_n \Phi)}{P(h^{(n)} \|Z\| > \lambda_n)} \leq \frac{
   \nu(B_{k\eps}^{-1}\Psi^{\delta} \cap C_{k\eps}^{-1}\Phi^{\delta})}{\nu(\{z: 
   \|z\| \geq 1\})},
 $$
 $$
\liminf_{n\to \infty} \frac{  P({H_{[kn\eps]}^{(n)}} Z_k \in \lambda_n
   \Psi, G_{[kn\eps]}^{(n)} Z_k \in \lambda_n \Phi)}{P(h^{(n)} \|Z\| > \lambda_n)} \geq \frac{  \nu(B_{k\eps}^{-1}\Psi^{-\delta} \cap C_{k\eps}^{-1} \Phi^{-\delta})}{\nu(\{z: \|z\| \geq 1\})}.$$

 % \nu (\overline{B_{k\eps}^{-1}\Gamma^{\delta}})
 
\end{lemma}

    \begin{equation}\label{eq:asym_equivalence}
        {P(S_0^n \in \lambda_n \Gamma)} \sim \frac{n P(h^{(n)} \|Z\| >
          \lambda_n)}{\nu(\{z: \|z\|\geq 1\})}  \int_{-\infty}^\infty
        \nu (B_y^{-1}\Gamma) \, dy. 
    \end{equation}

\begin{lemma}\label{lemma:cluster_equi_pt}
Suppose that $\Gamma$  satisfies Assumption \ref{ass:tail_measure}. Let
$0<\delta<\delta_0$ and denote 
\[
  D^{M,\pm\delta}(\theta_1,\theta_2) = \bigl\{(y, z): y \in [-M, M],\, B_{y-t}z
  \in \Psi^{\pm \delta} \ \text{for all} \
  t \in [-\theta_2, \theta_1]\cap \bbq\bigr\}.
\] 
% \problem{Maybe we need to remove the $\delta$ to achieve a matching bound? Do we need a certain continuity of the set?  Need to break it into upper and lower bounds, when to take $\delta \to 0$ and use continuity of set?}
Then for any $M^\prime >M$, $   0<\theta_1^\prime<\theta_1, \,
0<\theta_2^\prime<\theta_2$ we have 
\begin{equation} \label{e:connect.N2}
  \limsup_{n\to \infty} P\bigl(\cE_n^M(\Psi) \big|E_0(n,\Gamma)\bigr)
    - P\bigl(N_n^2\bigl(D^{M^\prime,\delta}(\theta_1^\prime, \theta_2^\prime\bigr)
    \geq 1\big|E_0(n,\Gamma)\bigr) \leq 0 
  \end{equation}
and for any $0<M^\prime<M$, $   \theta_1^\prime>\theta_1, \,
\theta_2^\prime>\theta_2$ we have
\begin{equation} \label{e:connect.N2a}
  \liminf_{n\to \infty} P\bigl(\cE_n^M(\Psi) \big|E_0(n,\Gamma)\bigr)
    - P\bigl(N_n^2\bigl(D^{M^\prime,-\delta}(\theta_1^\prime, \theta_2^\prime\bigr)
    \geq 1\big|E_0(n,\Gamma)\bigr) \geq 0.  
  \end{equation}
\end{lemma}

\begin{proof}
Fix   $M^\prime >M$, $   0<\theta_1^\prime<\theta_1, \,
0<\theta_2^\prime<\theta_2$ and let 
$M^{\prime\prime}\in (M,M^\prime)$. 
Let $\eps>0$ be a small number. For $-M/\eps\leq k\leq M/\eps$ let $J_k^{n, \eps} =
        [kn\eps, (k+1) n\eps) \cap \ZZ$ and denote $I_n^{M,\eps} =
        \bigcup_{-M/\eps\leq k\leq M/\eps} J_k^{n, \eps} $.  It is
        easy to check that
$$
[-(M-\eps)n,(M-\eps)n]\subseteq        I_n^{M,\eps} \subseteq
[-(M+\eps)n,(M+\eps)n].
$$
Let $\cE_n^{M,\eps}(\Psi)$ be defined identically to $\cE_n^M(\Psi)$
but with  $I_n^{M,\eps}$ replacing $I_n^M$. Let $\delta>0$ be a small
number. We conclude that for $0<\eps<\min(M^{\prime\prime} -M,
\theta_1-\theta_1^\prime, \theta_2-\theta_2^\prime)$ 
and all $n$ large
enough we have
\begin{align*}
 &P\bigl(\cE_n^M(\Psi) \big|E_0(n,\Gamma)\bigr) \leq 
  P\bigl(\cE_n^{M^\prime,\eps}(\Psi) \big|E_0(n,\Gamma)\bigr) \\
  \leq &P\bigl( \text{there is} \ -M^{\prime\prime}/\eps \leq k\leq
                    M^{\prime\prime}/\eps \ \text{and} \ i\in J_k^{n,
   \eps} \ \text{with}  \\
                  & \hskip 0.2in H^{(n)}_{[kn\eps]-j} Z_i\in \lambda_n \Psi^{\delta/4} \
                   \text{for all} \  j\in [-n\theta_2,n\theta_1] \big|E_0(n,\Gamma)\bigr)\\
\leq &P\bigl( \text{there is} \ -M^{\prime\prime}/\eps \leq k\leq
                    M^{\prime\prime}/\eps \ \text{and} \ i\in J_k^{n,
   \eps} \ \text{with}  \ H^{(n)}_{[(k-l)n\eps]} Z_i\in \lambda_n \Psi^{\delta/4} \\
                  & \hskip 0.2in 
                   \text{for all} \   -\theta_2^\prime/\eps\leq l\leq \theta_1^\prime/\eps \big|E_0(n,\Gamma)\bigr)
                    + P\bigl( R_1^{n, \eps}\big|E_0(n,\Gamma)\bigr)\\
\leq &P\bigl( \text{there is} \ -M^{\prime\prime}/\eps \leq k\leq
                    M^{\prime\prime}/\eps \ \text{and} \ i\in J_k^{n,
   \eps} \ \text{with}  \ B_{(k-l)\eps}h_m Z_i\in \lambda_n \Psi^{\delta/2} \\
                  & \hskip 0.2in  
                   \text{for all} \   -\theta_2^\prime/\eps\leq l\leq \theta_1^\prime/\eps \big|E_0(n,\Gamma)\bigr)
                    + P\bigl( R_1^{n,
                    \eps}\big|E_0(n,\Gamma)\bigr)+P\bigl( R_2^{n,
                    \eps}\big|E_0(n,\Gamma)\bigr) \\
\leq &P\Bigl( N_n^2\bigl(D_{\eps,n}^{M^{\prime\prime},\delta}(\theta_1^\prime,\theta_2^\prime)\bigr)\geq 1\big| E_0(n,\Gamma)\Bigr)
 + P\bigl( R_1^{n,
                    \eps}\big|E_0(n,\Gamma)\bigr)+P\bigl( R_2^{n,
       \eps}\big|E_0(n,\Gamma)\bigr) \\
  \leq &P\Bigl( N_n^2\bigl(D^{M^{\prime},\delta}(\theta_1^\prime,\theta_2^\prime)\bigr)\geq 1\big| E_0(n,\Gamma)\Bigr)
 + \sum_{l=1}^3 P\bigl( R_l^{n,
                    \eps}\big|E_0(n,\Gamma)\bigr), 
\end{align*}
where 
\begin{align*}
D_{\eps,n}^{M^\prime,\delta}(\theta_1^\prime,\theta_2^\prime) = &\bigl\{ (y,z):\,\text{there is} \ -M^\prime/\eps \leq k\leq
                   M^\prime/\eps \ \text{and} \; y\in J_k^{n,   \eps} /n \\
                   \
 & \text{with} \ B_{y-t}^\eps z\in \Psi^{\delta/2} \ \text{for all} \
  -\theta_2^\prime\leq t\leq \theta_1^\prime\bigr\}, 
\end{align*}
and the ``remainder'' events are defined by 
\begin{align*}
  R_1^{n, \eps}=& \bigl\{ \text{there is} \ -M^\prime/\eps \leq k\leq
                   M^\prime/\eps \ \text{and} \ i\in J_k^{n,
   \eps}  \\
   &\text{with} \ \| H_{[kn\eps]-j}-H_{i-j}\|\| Z_i\|
                          >\lambda_n\delta/4 \ \text{for some} \ j\in
                           [-n\theta_2,n\theta_1]\bigr\}, \\
 R_2^{n, \eps}=& \bigl\{ \text{there is} \   -M^\prime/\eps \leq k\leq
                   M^\prime/\eps \ \text{and} \ i\in J_k^{n,
   \eps} \\
&\text{with}  \ \|H^{(n)}_{[(k+l)n\eps]} -B_{(k+l)\eps}
                  h^{(n)}\|\|Z_i\|>\delta/4 \ 
                     \text{for some} \ -\theta_1^\prime/\eps\leq
                    l\leq\theta_2^\prime/\eps \bigr\}, \\
R_3^{n, \eps}=& \bigl\{ \text{there is} \ y\in [-M^{\prime},M^{\prime}] \
  \text{with} \ \| B_{y-t}-B_{y-t}^\eps\|h_n \|
                Z_{[yn]}\|>\lambda_n\delta/2 \\
   &\text{for some} \
                -\theta_2^\prime\leq t\leq \theta_1^\prime \bigr\}. 
   \end{align*}
We see that the claim \eqref{e:connect.N2} will follow once we prove
that 
   \begin{equation} \label{e:remainders.small}
 \lim_{\eps\to 0} \limsup_{n\to\infty} P\bigl( R_d^{n,
                    \eps}\big|E_0(n,\Gamma)\bigr) =0 \ \text{for} \ d=1,2,3.
                \end{equation}

For $d=1$ we use \eqref{eq:coeff_group_err} to note that for a fixed
$\eps>0$ and all $n$ large enough, for all relevant $k$ in $j$ we have 
  $\| H_{[kn\eps]-j}-H_{i-j}\| \leq f(\eps) h_n $ for  some $f(\eps)
  \to 0$ as $\eps \to 0$.  Therefore, by \eqref{eq:asym_equivalence}, 
 \begin{align*}
\limsup_{n\to\infty} P\bigl( R_1^{n,
                    \eps}\big|E_0(n,\Gamma)\bigr) 
\leq& 2M^\prime (\theta_1+\theta_2)\limsup_{n\to\infty}  \frac{nP\bigl(f(\eps) h^{(n)} \|Z\|
    > \lambda_n \delta/4\bigr)}{P\bigl( E_0(n,\Gamma)\bigr)} \\
=&O((f(\eps))^\alpha),
  \end{align*}
and   \eqref{e:remainders.small} with $d=1$ follows. 

Next, the continuity of $B_y$ as a function of $y$ means that
$$
\sup_{y\in [-M^\prime,M^\prime], \, t \in [-\theta_2^\prime,
  \theta_1^\prime]} \|B_{y-t} - B_{y - t}^\eps\| \leq g(\eps)
$$
 for some $g(\eps)\to 0$ as $\eps\to 0$. Therefore, repeating the
 argument in the case $d=1$ we see that for a fixed $\eps>0$,
 $$
 \limsup_{n\to\infty} P\bigl( R_3^{n,
                    \eps}\big|E_0(n,\Gamma)\bigr) = O((g(\eps))^\alpha),
                  $$
giving us \eqref{e:remainders.small} with $d=3$, and so it remains to
prove \eqref{e:remainders.small} for $d=2$. Consider all relevant
choices of $k$ and $l$. Taking the case $k+l\leq 0$ first, we write 
\begin{align*}
     & \sup_{k,l:\, k+l\leq 0}\left\|H_{[(k+l)n\eps]}^{(n)}/h^{(n)}-
       B_{(k+l)\eps}\right\| \\
&\leq  \sup_{k,l:\, k+l\leq 0} \left\|H_0^{n-[(k+l)n\eps]} /h^{(n)}- (1-(k+l)\eps)^{1+\beta}B_+ \right\|  \\
     & +  \sup_{k,l:\, k+l\leq 0} \left\|H_0^{[(-(k+l) n\eps)]}/h^{(n)}
     - (-(k+l)\eps)^{1+\beta} B_+ \right\| =   o(1)
\end{align*}
as we have seen before (in the proof of Lemma
\ref{lemma:asym_lim_meas}). Since the case $k+l>0$ is similar, 
 \eqref{e:remainders.small} for $d=2$ follows using the same arguments
 as for the case $d=1$. 

 This proves \eqref{e:connect.N2}, and \eqref{e:connect.N2a} can be
 proved in the same way. 
\end{proof}

\end{document}
